\section{Lattice counts with a prescribed determinant}
\label{lat:section}

We prove a uniform bound for integer matrices whose rows belong to a fixed
lattice and whose determinant is a prescribed nonzero integer.  The lattice
may have very different scales in different directions, so the proof keeps
track of three separate column bounds.  We also record the plane counts
needed when the determinant is zero.

All lengths in this section are Euclidean.  The determinant of a lattice
means its covolume in its real span.  For a lattice $L$ of positive rank $s$, choose
$v_1,\ldots,v_s$ successively, with $v_i$ a shortest lattice vector outside
$\operatorname{span}_{\mathbb R}(v_1,\ldots,v_{i-1})$, and put
$\lambda_i=|v_i|$.  In particular,
$0<\lambda_1\le\cdots\le\lambda_s$.  The notation $A\ll B$ means
$A\le CB$ for an absolute constant $C$, unless dependence is indicated.

The classical geometry-of-numbers antecedent is Minkowski's second
theorem on successive minima; see Henk~\cite[Theorem 1.3]{Henk2002}.
Only the weaker fixed-dimensional Euclidean product bound below is
needed here, and we include its proof.

\begin{lemma}[Successive lengths]\label{lat:successive}
For a lattice $L$ of positive rank $s$ and the lengths just defined,
\[
 \det L\le\prod_{i=1}^s\lambda_i\le s^{s/2}\det L.
\]
Consequently, the lattice generated by $v_1,\ldots,v_s$ has index at most
$s^{s/2}$ in $L$.
\end{lemma}

\begin{proof}
The determinant of the lattice generated by the $v_i$ is an integer
multiple of $\det L$ and is at most $\prod_i\lambda_i$.
For the reverse bound, choose orthonormal coordinates whose first $j$
directions span $v_1,\ldots,v_j$ for each $j$.  Let
\[
 Q=\{(a_1,\ldots,a_s): |a_i|<\lambda_i/\sqrt{s}\text{ for every }i\}.
\]
If a nonzero lattice vector in $Q$ has last nonzero coordinate $j$, it
lies outside the preceding span and has length strictly less than
$\lambda_j$, a contradiction.  The translates of $Q/2$ by $L$ are
therefore disjoint.  Integrating their indicator functions over a
fundamental domain of $L$ gives
$\operatorname{vol}(Q/2)\le\det L$.
Since $\operatorname{vol}(Q/2)=s^{-s/2}\prod_i\lambda_i$, the result follows.
\end{proof}

\begin{lemma}[Points in a plane lattice]\label{lat:plane-count}
Let $L$ be a rank-two Euclidean lattice with successive lengths
$\lambda_1\le\lambda_2$.  If $R\ge\lambda_2$, any translate of $L$ has
at most $C R^2/\det L$ points in any ball of radius $R$.
If $R<\lambda_2$, the points of $L$ in the ball of radius $R$ centered
at zero lie in a line and number at most $1+2R/\lambda_1$.
\end{lemma}

\begin{proof}
For the first assertion let $L_0=\mathbb Zv_1+\mathbb Zv_2$.
A half-open fundamental parallelogram for $L_0$ has diameter at most
$\lambda_1+\lambda_2\le2R$.  For each coset of $L_0$, attach such a
parallelogram to every point in the ball.  The parallelograms are
disjoint and lie in a disk of radius at most $3R$ in the affine span.
Thus each coset contributes at most $9\pi R^2/\det L_0$ points.
Summing over $[L:L_0]$ cosets proves the first assertion.
For the second, every lattice vector outside $\mathbb Rv_1$ has length
at least $\lambda_2$.  The lattice on $\mathbb Rv_1$ is
$\mathbb Zv_1$, since a shorter generator would contradict the choice
of $v_1$.  Counting its multiples proves the bound.
\end{proof}

\begin{lemma}[Covolumes of integral planes]\label{lat:plane-covolume}
Let $y\in\mathbb Z^3$ be primitive, meaning that its coordinates have
greatest common divisor one.  Then
$L_y=\mathbb Z^3\cap y^\perp$ has determinant $|y|$.
More generally, if $\Lambda\subseteq\mathbb Z^3$ is a full-rank lattice
of index $I$ and $y\cdot\Lambda=g\mathbb Z$ with $g>0$, then
\[
 \det(\Lambda\cap y^\perp)=\frac{I|y|}{g}.
\]
\end{lemma}

\begin{proof}
Choose $u\in\Lambda$ with $y\cdot u=g$.  The map
$v\mapsto y\cdot v/g$ maps $\Lambda$ onto $\mathbb Z$, has kernel
$\Lambda\cap y^\perp$, and splits by $1\mapsto u$.  Hence a basis of
the kernel together with $u$ is a basis of $\Lambda$.
The height of $u$ above $y^\perp$ is $g/|y|$, which proves the formula.
For $\Lambda=\mathbb Z^3$, primitivity gives $g=1$.
\end{proof}

We next count matrices with unequal column bounds.  The difficult case
occurs when the third-column radius reaches only one direction of its
integral plane; we count those planes through their shortest vectors.

\begin{lemma}[Unequal column bounds]\label{lat:anisotropic}
Let $R_1\ge R_2\ge R_3\ge1$ and let $t\in\mathbb Z\setminus\{0\}$.
The number of triples $u_1,u_2,u_3\in\mathbb Z^3$ satisfying
$|u_j|\le R_j$ and $\det(u_1,u_2,u_3)=t$ is at most
\[
 C(R_1R_2R_3)^2\bigl(\log(2R_1)\bigr)^2.
\]
The constant is independent of $t$ and of the three radii.
\end{lemma}

\begin{proof}
The columns $u_2,u_3$ are independent.  Choose a primitive normal $y$
to their plane, fixing its sign by requiring its first nonzero coordinate
to be positive.  Then $u_2\times u_3=a y$ for some nonzero integer $a$,
so $|y|\le R_2R_3$.  If $\mu_1,\mu_2$ are the successive lengths of
$L_y$, the columns imply $\mu_1\le R_3$ and $\mu_2\le R_2$.

For fixed $u_2,u_3$, the determinant equation is $u_1\cdot y=t/a$.
Its integral solutions are either absent or form a translate of $L_y$.
By Lemmas~\ref{lat:plane-count} and~\ref{lat:plane-covolume}, at most
$C R_1^2/|y|$ satisfy the first-column bound; the required hypothesis is
$\mu_2\le R_2\le R_1$.  For a fixed normal $y$, there are likewise at
most $C R_2^2/|y|$ choices for $u_2$.

First suppose $\mu_2\le R_3$.  There are at most
$C R_3^2/|y|$ choices for $u_3$.  The resulting bound, summed over
normals, is
\[
 C(R_1R_2R_3)^2
 \sum_{\substack{y\in\mathbb Z^3\\0<|y|\le R_2R_3}}|y|^{-3}
 \ll (R_1R_2R_3)^2\log(2R_2R_3).
\]
The last estimate follows by splitting the integer vectors into dyadic
shells; a shell of radius $Y\ge1$ has $O(Y^3)$ vectors.

It remains to treat $\mu_2>R_3$.  Group the normals by powers of two
$U,Z\ge1$ with
\[
 U\le\mu_1<2U,\qquad Z\le\mu_2<2Z.
\]
Lemma~\ref{lat:successive} gives $|y|\asymp UZ$.
There are at most $C U^4Z^2$ normals in such a group.  Indeed, choose
a shortest nonzero vector $z\in L_y$.  It is primitive in $\mathbb Z^3$,
because division by a common divisor would produce a shorter vector
in $L_y$.  Its length is in $[U,2U)$, giving $O(U^3)$ choices.
For each such $z$, the possible normals lie in $L_z$ and have length
less than $4UZ$.  If $\nu_1,\nu_2$ are the successive lengths of
$L_z$, integrality gives $\nu_1\ge1$, and the preceding lemmas give
\[
 \nu_2\le\nu_1\nu_2\le2|z|<4U\le4UZ.
\]
Consequently Lemma~\ref{lat:plane-count}, applied at radius $4UZ$,
bounds their number by
$C(UZ)^2/|z|\le C UZ^2$.  This proves the claimed normal count.

For each of these normals, all eligible third columns lie on the line
of a shortest vector.  Lemma~\ref{lat:plane-count} gives at most
$1+2R_3/\mu_1\ll R_3/U$ choices, since $U\le\mu_1\le R_3$.
Thus the contribution of one dyadic group is at most
\[
 C U^4Z^2\,
 \frac{R_1^2}{UZ}\,
 \frac{R_2^2}{UZ}\,
 \frac{R_3}{U}
 =C R_1^2R_2^2R_3U
 \le C(R_1R_2R_3)^2.
\]
There are $O(\log(2R_3)\log(2R_2))$ groups, since
$1\le U\le R_3$ and $1\le Z\le R_2$.  Combining both cases proves
the assertion.
\end{proof}

\begin{proposition}[A fixed nonzero determinant]\label{lat:fixed-det}
Let $\Lambda\subseteq\mathbb Z^3$ have index $I$, let $X\ge2$, and
suppose its third successive length satisfies $\lambda_3\le X$.
For every $t\in\mathbb Z\setminus\{0\}$, the number of $3$ by $3$
matrices with rows in $\Lambda$, row lengths at most $X$, and
determinant $t$ is at most
\[
 C\bigl(\log(2X)\bigr)^2\frac{X^6}{I^2}.
\]
The constant is absolute.
\end{proposition}

\begin{proof}
Let $V$ be the matrix whose rows are the successive shortest vectors
$v_1,v_2,v_3$, and let $\Lambda_0$ be their integer span.  By
Lemma~\ref{lat:successive},
\[
 m=[\Lambda:\Lambda_0]=\frac{|\det V|}{I}\le3^{3/2}.
\]
Since the finite group $\Lambda/\Lambda_0$ has order $m$, each
$u\in\Lambda$ has coordinates in $m^{-1}\mathbb Z$ relative to this
basis.  If $u=\sum_j\alpha_jv_j$ and $|u|\le X$, Cramer's rule gives
\[
 |\alpha_j|
 \le\frac{X\prod_{i\ne j}\lambda_i}{|\det V|}
 \le 3^{3/2}\frac{X}{\lambda_j}.
\]
Thus the map $A\mapsto Q=mAV^{-1}$ is injective and takes the matrices
under consideration to integral matrices.  Their common determinant is
$m^3t/\det V\ne0$.  If this number is not an integer, there are no
such matrices; otherwise Lemma~\ref{lat:anisotropic} applies.

The $j$th column of $Q$ has length at most
$R_j=C_0X/\lambda_j$, where $C_0$ is an absolute constant chosen
sufficiently large.  These radii satisfy $R_1\ge R_2\ge R_3\ge1$.
Moreover, integrality of $\Lambda$ gives $\lambda_j\ge1$, and
Lemma~\ref{lat:successive} gives
\[
 R_j\le C_0X,\qquad
 R_1R_2R_3=\frac{C_0^3X^3}{\lambda_1\lambda_2\lambda_3}
 \le\frac{C_0^3X^3}{I}.
\]
The bound in Lemma~\ref{lat:anisotropic} is therefore the claimed one.
\end{proof}
